% ===================================================================
% Section: Priklad navrhu implementacie)
% ===================================================================

\section{Priklad navrhu implementacie}
\label{sec:intro_example_control}

Začatie tejto práce bolo zvolené formou opisného úvodného príkladu regulácie,
konkrétne štandardného a všeobecne známeho PID regulátora.
Predpokladajme, že máme k dispozícii identifikovaný model systému,
prípadne uvažujeme priamu implementáciu algoritmu na reálnom systéme.
V oboch prípadoch budeme riadený objekt ďalej označovať ako $G_s(s)$.

% -------------------------------------------------------------------
% Subsection: PID regulácia v uzavretej spätnej väzbe 
% -------------------------------------------------------------------
\subsection{PID regulácia v uzavretej spätnej väzbe}

Prenosová funkcia PID regulátora v paralelnej forme je definovaná ako
\begin{equation}
    C(s) = K_p + \frac{K_i}{s} + K_d s,
    \label{eq:intro_pid_parallel}
\end{equation}
kde $K_p$ predstavuje proporcionálny zisk, $K_i$ integračný zisk
a $K_d$ derivačný zisk regulátora.

Na reguláciu budeme uvažovať štandardnú zápornú spätnú väzbu.
Regulačná odchýlka vzniká ako rozdiel medzi požadovanou hodnotou a výstupom systému,
\begin{equation}
    e(t) = w(t) - y(t),
    \label{eq:intro_error_cont}
\end{equation}
pričom $w(t)$ je žiadaná hodnota (setpoint) a $y(t)$ je výstup regulovanej sústavy.
Výstup regulátora predstavuje akčný zásah $u(t)$, ktorý vstupuje do systému $G_s(s)$.

Pre diskrétnu implementáciu s periódou vzorkovania $T_s$ budeme používať zápis
\begin{equation}
    e(k) = w(k) - y(k),
    \label{eq:intro_error_disc}
\end{equation}
kde index $k$ označuje diskrétny časový krok.

Bloková schéma uzavretého regulačného obvodu má nasledujúci tvar:
\begin{center}
\vbox{
\makebox[\textwidth][c]{%
    \includegraphics[width=1.0\textwidth]{URO_scheme.pdf}
}
\figcaption{Bloková schéma uzavretého regulačného obvodu so zápornou spätnou väzbou a PID regulátorom}
\label{fig:uro_pid_block_scheme}
}
\end{center}

% -------------------------------------------------------------------
% Subsection: Prenosové funkcie ORO/URO a zákon regulácie 
% -------------------------------------------------------------------
\subsection{Prenosové funkcie ORO a URO a zákon regulácie}

Predpokladajme, že regulovaný systém má prenosovú funkciu s čitateľom nultého rádu a menovateľom druhého rádu,
\begin{equation}
    G_s(s) = \frac{b_0}{s^2 + a_1 s + a_0}.
    \label{eq:intro_plant_2nd_order}
\end{equation}

Prenosová funkcia otvoreného regulačného obvodu (ORO) je daná súčinom systému a regulátora,
\begin{equation}
    G_{\mathrm{ORO}}(s) = G_s(s)\,C(s).
    \label{eq:intro_oro_def}
\end{equation}

Po dosadení \eqref{eq:intro_pid_parallel} a \eqref{eq:intro_plant_2nd_order} dostaneme
\begin{align}
    G_{\mathrm{ORO}}(s)
    &= \frac{b_0}{s^2 + a_1 s + a_0}
       \left(K_p + \frac{K_i}{s} + K_d s\right) \nonumber \\
    &= \frac{b_0\left(K_d s^2 + K_p s + K_i\right)}{s\left(s^2 + a_1 s + a_0\right)}.
    \label{eq:intro_oro_result}
\end{align}

Prenosová funkcia uzavretého regulačného obvodu (URO) pri jednotkovej spätnej väzbe je
\begin{equation}
    G_{\mathrm{URO}}(s) = \frac{G_{\mathrm{ORO}}(s)}{1 + G_{\mathrm{ORO}}(s)}.
    \label{eq:intro_uro_def}
\end{equation}
Uvedený tvar vychádza zo štandardného opisu uzavretej spätnej väzby
v klasickej teórii riadenia~\cite{nise2015}.

Po dosadení \eqref{eq:intro_oro_result} platí
\begin{equation}
    G_{\mathrm{URO}}(s)
    =
    \frac{b_0\left(K_d s^2 + K_p s + K_i\right)}
    {s\left(s^2 + a_1 s + a_0\right) + b_0\left(K_d s^2 + K_p s + K_i\right)}.
    \label{eq:intro_uro_result}
\end{equation}

Póly ORO sú dané koreňmi menovateľa prenosu $G_{\mathrm{ORO}}(s)$, t.\,j.
\begin{equation}
    s\left(s^2 + a_1 s + a_0\right) = 0,
    \label{eq:intro_oro_poles_eq}
\end{equation}
odkiaľ priamo vyplýva jeden pól v $s_1 = 0$ a zvyšné dva póly sú koreňmi kvadratickej rovnice
\begin{equation}
    s^2 + a_1 s + a_0 = 0.
    \label{eq:intro_plant_poles_eq}
\end{equation}

Póly URO sú určené koreňmi charakteristického polynómu (menovateľa) uzavretého obvodu,
\begin{equation}
    s\left(s^2 + a_1 s + a_0\right) + b_0\left(K_d s^2 + K_p s + K_i\right) = 0,
    \label{eq:intro_uro_poles_eq}
\end{equation}
pričom po roznásobení dostaneme explicitný tvar tretieho rádu
\begin{equation}
    s^3 + \left(a_1 + b_0 K_d\right)s^2 + \left(a_0 + b_0 K_p\right)s + b_0 K_i = 0.
    \label{eq:intro_uro_char_poly}
\end{equation}

Stabilita je v oboch prípadoch kľúčová, avšak význam sa líši. V prípade ORO sa stabilita posudzuje najmä z hľadiska robustnosti a bezpečnosti návrhu regulátora. V prípade URO je stabilita nevyhnutná podmienka správnej funkcie regulácie. Pre spojité systémy stabilita znamená, že všetky póly $G_{\mathrm{URO}}(s)$ majú zápornú reálnu časť \cite{ogata2010}. V opačnom prípade výstup $y(t)$ diverguje, regulačná odchýlka neklesá a systém môže byť v reálnej prevádzke nebezpečný.


% -----------------------------
% Control law and discretization
% -----------------------------
\paragraph{Zákon riadenia}

V časovej oblasti je PID zákon regulácie možné zapísať ako
\begin{equation}
    u(t) = K_p e(t) + K_i \int_0^t e(\tau)\,d\tau + K_d \frac{d e(t)}{dt}.
    \label{eq:intro_pid_law_cont}
\end{equation}

Pre praktickú implementáciu v diskrétnom čase budeme uvažovať základnú aproximáciu:
\begin{align}
    I(k) &= I(k-1) + T_s\,e(k), \label{eq:intro_int_update} \\
    D(k) &= \frac{e(k) - e(k-1)}{T_s}, \label{eq:intro_der_update} \\
    u(k) &= K_p e(k) + K_i I(k) + K_d D(k).
    \label{eq:intro_pid_law_disc}
\end{align}

Z dôvodu fyzikálnych obmedzení akčného zásahu sa v praxi uplatňuje saturácia
\begin{equation}
    u(k) = \min\bigl(\max(u(k),u_{\min}),u_{\max}\bigr).
    \label{eq:intro_saturation}
\end{equation}

% -----------------------------
% Pseudocode
% -----------------------------

\paragraph{Pseudokód PID regulátora}

Nižšie je uvedený ilustračný pseudokód výpočtu PID regulátora v diskrétnom čase
s vnútornými stavmi integrátora a pamäťou predchádzajúcej chyby.

\begin{lstlisting}[language=Matlab,
caption={Pseudokód diskrétnej implementácie PID regulátora (ilustračný príklad)},
label={lst:intro_pid_pseudocode}]
% Inputs: w (setpoint), y (output), Ts
% States: I (integrator state), e_prev (previous error)
% Params: Kp, Ki, Kd, u_min, u_max

e = w - y;

% Integrator update
I = I + Ts * e;

% Derivative term
D = (e - e_prev) / Ts;

% PID law
u = Kp * e + Ki * I + Kd * D;

% Saturation
u = min(max(u, u_min), u_max);

% Memory update
e_prev = e;
\end{lstlisting}